\documentclass[12pt]{article}

% Idioma y codificación
\usepackage[spanish, es-tabla]{babel}       %es-tabla para que se titule "Tabla"
\usepackage[utf8]{inputenc}

% Márgenes
\usepackage[a4paper,top=3cm,bottom=2.5cm,left=3cm,right=3cm]{geometry}

% Comentarios de bloque
\usepackage{verbatim}

% Paquetes de links
\usepackage[hidelinks]{hyperref}    % Permite enlaces
\usepackage{url}                    % redirecciona a la web

% Más opciones para enumeraciones
\usepackage{enumitem}

% Personalizar la portada
\usepackage{titling}


% Paquetes de tablas
\usepackage{multirow}


%------------------------------------------------------------------------

%Paquetes de figuras
\usepackage{caption}
\usepackage{subcaption} % Figuras al lado de otras
\usepackage{float}      % Poner figuras en el sitio indicado H.


% Paquetes de imágenes
\usepackage{graphicx}       % Paquete para añadir imágenes
\usepackage{transparent}    % Para manejar la opacidad de las figuras

% Paquete para usar colores
\usepackage[dvipsnames]{xcolor}
\usepackage{pagecolor}      % Para cambiar el color de la página

% Habilita tamaños de fuente mayores
\usepackage{fix-cm}


%------------------------------------------------------------------------

% Paquetes de matemáticas
\usepackage{mathtools, amsfonts, amssymb, mathrsfs}
\usepackage[makeroom]{cancel}     % Simplificar tachando
\usepackage{polynom}    % Divisiones y Ruffini
\usepackage{units} % Para poner fracciones diagonales con \nicefrac

\usepackage{pgfplots}   %Representar funciones
\pgfplotsset{compat=1.18}  % Versión 1.18

\usepackage{tikz-cd}    % Para usar diagramas de composiciones
\usetikzlibrary{calc}   % Para usar cálculo de coordenadas en tikz

%Definición de teoremas, etc.
\usepackage{amsthm}
%\swapnumbers   % Intercambia la posición del texto y de la numeración

\theoremstyle{plain}

\makeatletter
\@ifclassloaded{article}{
  \newtheorem{teo}{Teorema}[section]
}{
  \newtheorem{teo}{Teorema}[chapter]  % Se resetea en cada chapter
}
\makeatother

\newtheorem{coro}{Corolario}[teo]           % Se resetea en cada teorema
\newtheorem{prop}[teo]{Proposición}         % Usa el mismo contador que teorema
\newtheorem{lema}[teo]{Lema}                % Usa el mismo contador que teorema

\theoremstyle{remark}
\newtheorem*{observacion}{Observación}

\theoremstyle{definition}

\makeatletter
\@ifclassloaded{article}{
  \newtheorem{definicion}{Definición} [section]     % Se resetea en cada chapter
}{
  \newtheorem{definicion}{Definición} [chapter]     % Se resetea en cada chapter
}
\makeatother

\newtheorem*{notacion}{Notación}
\newtheorem*{ejemplo}{Ejemplo}
\newtheorem*{ejercicio*}{Ejercicio}             % No numerado
\newtheorem{ejercicio}{Ejercicio} [section]     % Se resetea en cada section


% Modificar el formato de la numeración del teorema "ejercicio"
\renewcommand{\theejercicio}{%
  \ifnum\value{section}=0 % Si no se ha iniciado ninguna sección
    \arabic{ejercicio}% Solo mostrar el número de ejercicio
  \else
    \thesection.\arabic{ejercicio}% Mostrar número de sección y número de ejercicio
  \fi
}


% \renewcommand\qedsymbol{$\blacksquare$}         % Cambiar símbolo QED
%------------------------------------------------------------------------

% Paquetes para encabezados
\usepackage{fancyhdr}
\pagestyle{fancy}
\fancyhf{}

\newcommand{\helv}{ % Modificación tamaño de letra
\fontfamily{}\fontsize{12}{12}\selectfont}
\setlength{\headheight}{15pt} % Amplía el tamaño del índice


%\usepackage{lastpage}   % Referenciar última pag   \pageref{LastPage}
\fancyfoot[C]{\thepage}

%------------------------------------------------------------------------

% Conseguir que no ponga "Capítulo 1". Sino solo "1."
\makeatletter
\@ifclassloaded{book}{
  \renewcommand{\chaptermark}[1]{\markboth{\thechapter.\ #1}{}} % En el encabezado
    
  \renewcommand{\@makechapterhead}[1]{%
  \vspace*{50\p@}%
  {\parindent \z@ \raggedright \normalfont
    \ifnum \c@secnumdepth >\m@ne
      \huge\bfseries \thechapter.\hspace{1em}\ignorespaces
    \fi
    \interlinepenalty\@M
    \Huge \bfseries #1\par\nobreak
    \vskip 40\p@
  }}
}
\makeatother

%------------------------------------------------------------------------
% Paquetes de cógido
\usepackage{minted}
\renewcommand\listingscaption{Código fuente}

\usepackage{fancyvrb}
% Personaliza el tamaño de los números de línea
\renewcommand{\theFancyVerbLine}{\small\arabic{FancyVerbLine}}

% Estilo para C++
\newminted{cpp}{
    frame=lines,
    framesep=2mm,
    baselinestretch=1.2,
    linenos,
    escapeinside=||
}



\usepackage{listings} % Para incluir código desde un archivo

\renewcommand\lstlistingname{Código Fuente}
\renewcommand\lstlistlistingname{Índice de Códigos Fuente}

% Definir colores
\definecolor{vscodepurple}{rgb}{0.5,0,0.5}
\definecolor{vscodeblue}{rgb}{0,0,0.8}
\definecolor{vscodegreen}{rgb}{0,0.5,0}
\definecolor{vscodegray}{rgb}{0.5,0.5,0.5}
\definecolor{vscodebackground}{rgb}{0.97,0.97,0.97}
\definecolor{vscodelightgray}{rgb}{0.9,0.9,0.9}

% Configuración para el estilo de C similar a VSCode
\lstdefinestyle{vscode_C}{
  backgroundcolor=\color{vscodebackground},
  commentstyle=\color{vscodegreen},
  keywordstyle=\color{vscodeblue},
  numberstyle=\tiny\color{vscodegray},
  stringstyle=\color{vscodepurple},
  basicstyle=\scriptsize\ttfamily,
  breakatwhitespace=false,
  breaklines=true,
  captionpos=b,
  keepspaces=true,
  numbers=left,
  numbersep=5pt,
  showspaces=false,
  showstringspaces=false,
  showtabs=false,
  tabsize=2,
  frame=tb,
  framerule=0pt,
  aboveskip=10pt,
  belowskip=10pt,
  xleftmargin=10pt,
  xrightmargin=10pt,
  framexleftmargin=10pt,
  framexrightmargin=10pt,
  framesep=0pt,
  rulecolor=\color{vscodelightgray},
  backgroundcolor=\color{vscodebackground},
}

%------------------------------------------------------------------------

% Comandos definidos
\newcommand{\bb}[1]{\mathbb{#1}}
\newcommand{\cc}[1]{\mathcal{#1}}

% I prefer the slanted \leq
\let\oldleq\leq % save them in case they're every wanted
\let\oldgeq\geq
\renewcommand{\leq}{\leqslant}
\renewcommand{\geq}{\geqslant}

% Si y solo si
\newcommand{\sii}{\iff}

% Letras griegas
\newcommand{\eps}{\epsilon}
\newcommand{\veps}{\varepsilon}
\newcommand{\lm}{\lambda}

\newcommand{\ol}{\overline}
\newcommand{\ul}{\underline}
\newcommand{\wt}{\widetilde}
\newcommand{\wh}{\widehat}

\let\oldvec\vec
\renewcommand{\vec}{\overrightarrow}

% Derivadas parciales
\newcommand{\del}[2]{\frac{\partial #1}{\partial #2}}
\newcommand{\Del}[3]{\frac{\partial^{#1} #2}{\partial^{#1} #3}}
\newcommand{\deld}[2]{\dfrac{\partial #1}{\partial #2}}
\newcommand{\Deld}[3]{\dfrac{\partial^{#1} #2}{\partial^{#1} #3}}


\newcommand{\AstIg}{\stackrel{(\ast)}{=}}
\newcommand{\Hop}{\stackrel{L'H\hat{o}pital}{=}}

\newcommand{\red}[1]{{\color{red}#1}} % Para integrales, destacar los cambios.

% Método de integración
\newcommand{\MetInt}[2]{
    \left[\begin{array}{c}
        #1 \\ #2
    \end{array}\right]
}

% Declarar aplicaciones
% 1. Nombre aplicación
% 2. Dominio
% 3. Codominio
% 4. Variable
% 5. Imagen de la variable
\newcommand{\Func}[5]{
    \begin{equation*}
        \begin{array}{rrll}
            #1:& #2 & \longrightarrow & #3\\
               & #4 & \longmapsto & #5
        \end{array}
    \end{equation*}
}

%------------------------------------------------------------------------


\begin{document}

    % 1. Foto de fondo
    % 2. Título
    % 3. Encabezado Izquierdo
    % 4. Color de fondo
    % 5. Coord x del titulo
    % 6. Coord y del titulo
    % 7. Fecha

    
    % 1. Foto de fondo
% 2. Título
% 3. Encabezado Izquierdo
% 4. Color de fondo
% 5. Coord x del titulo
% 6. Coord y del titulo
% 7. Fecha

\newcommand{\portada}[7]{

    \portadaBase{#1}{#2}{#3}{#4}{#5}{#6}{#7}
    \portadaBook{#1}{#2}{#3}{#4}{#5}{#6}{#7}
}

\newcommand{\portadaExamen}[7]{

    \portadaBase{#1}{#2}{#3}{#4}{#5}{#6}{#7}
    \portadaArticle{#1}{#2}{#3}{#4}{#5}{#6}{#7}
}




\newcommand{\portadaBase}[7]{

    % Tiene la portada principal y la licencia Creative Commons
    
    % 1. Foto de fondo
    % 2. Título
    % 3. Encabezado Izquierdo
    % 4. Color de fondo
    % 5. Coord x del titulo
    % 6. Coord y del titulo
    % 7. Fecha
    
    
    \thispagestyle{empty}               % Sin encabezado ni pie de página
    \newgeometry{margin=0cm}        % Márgenes nulos para la primera página
    
    
    % Encabezado
    \fancyhead[L]{\helv #3}
    \fancyhead[R]{\helv \nouppercase{\leftmark}}
    
    
    \pagecolor{#4}        % Color de fondo para la portada
    
    \begin{figure}[p]
        \centering
        \transparent{0.3}           % Opacidad del 30% para la imagen
        
        \includegraphics[width=\paperwidth, keepaspectratio]{assets/#1}
    
        \begin{tikzpicture}[remember picture, overlay]
            \node[anchor=north west, text=white, opacity=1, font=\fontsize{60}{90}\selectfont\bfseries\sffamily, align=left] at (#5, #6) {#2};
            
            \node[anchor=south east, text=white, opacity=1, font=\fontsize{12}{18}\selectfont\sffamily, align=right] at (9.7, 3) {\textbf{\href{https://losdeldgiim.github.io/}{Los Del DGIIM}}};
            
            \node[anchor=south east, text=white, opacity=1, font=\fontsize{12}{15}\selectfont\sffamily, align=right] at (9.7, 1.8) {Doble Grado en Ingeniería Informática y Matemáticas\\Universidad de Granada};
        \end{tikzpicture}
    \end{figure}
    
    
    \restoregeometry        % Restaurar márgenes normales para las páginas subsiguientes
    \pagecolor{white}       % Restaurar el color de página
    
    
    \newpage
    \thispagestyle{empty}               % Sin encabezado ni pie de página
    \begin{tikzpicture}[remember picture, overlay]
        \node[anchor=south west, inner sep=3cm] at (current page.south west) {
            \begin{minipage}{0.5\paperwidth}
                \href{https://creativecommons.org/licenses/by-nc-nd/4.0/}{
                    \includegraphics[height=2cm]{assets/Licencia.png}
                }\vspace{1cm}\\
                Esta obra está bajo una
                \href{https://creativecommons.org/licenses/by-nc-nd/4.0/}{
                    Licencia Creative Commons Atribución-NoComercial-SinDerivadas 4.0 Internacional (CC BY-NC-ND 4.0).
                }\\
    
                Eres libre de compartir y redistribuir el contenido de esta obra en cualquier medio o formato, siempre y cuando des el crédito adecuado a los autores originales y no persigas fines comerciales. 
            \end{minipage}
        };
    \end{tikzpicture}
    
    
    
    % 1. Foto de fondo
    % 2. Título
    % 3. Encabezado Izquierdo
    % 4. Color de fondo
    % 5. Coord x del titulo
    % 6. Coord y del titulo
    % 7. Fecha


}


\newcommand{\portadaBook}[7]{

    % 1. Foto de fondo
    % 2. Título
    % 3. Encabezado Izquierdo
    % 4. Color de fondo
    % 5. Coord x del titulo
    % 6. Coord y del titulo
    % 7. Fecha

    % Personaliza el formato del título
    \pretitle{\begin{center}\bfseries\fontsize{42}{56}\selectfont}
    \posttitle{\par\end{center}\vspace{2em}}
    
    % Personaliza el formato del autor
    \preauthor{\begin{center}\Large}
    \postauthor{\par\end{center}\vfill}
    
    % Personaliza el formato de la fecha
    \predate{\begin{center}\huge}
    \postdate{\par\end{center}\vspace{2em}}
    
    \title{#2}
    \author{\href{https://losdeldgiim.github.io/}{Los Del DGIIM}}
    \date{Granada, #7}
    \maketitle
    
    \tableofcontents
}




\newcommand{\portadaArticle}[7]{

    % 1. Foto de fondo
    % 2. Título
    % 3. Encabezado Izquierdo
    % 4. Color de fondo
    % 5. Coord x del titulo
    % 6. Coord y del titulo
    % 7. Fecha

    % Personaliza el formato del título
    \pretitle{\begin{center}\bfseries\fontsize{42}{56}\selectfont}
    \posttitle{\par\end{center}\vspace{2em}}
    
    % Personaliza el formato del autor
    \preauthor{\begin{center}\Large}
    \postauthor{\par\end{center}\vspace{3em}}
    
    % Personaliza el formato de la fecha
    \predate{\begin{center}\huge}
    \postdate{\par\end{center}\vspace{5em}}
    
    \title{#2}
    \author{\href{https://losdeldgiim.github.io/}{Los Del DGIIM}}
    \date{Granada, #7}
    \thispagestyle{empty}               % Sin encabezado ni pie de página
    \maketitle
    \vfill
}
    \portadaExamen{ffccA4.jpg}{Ecuaciones\\Diferenciales I\\Examen III}{Ecuaciones Diferenciales I. Examen III}{MidnightBlue}{-8}{28}{2024-2025}{Arturo Olivares Martos}

    \begin{description}
        \item[Asignatura] Ecuaciones Diferenciales I
        \item[Curso Académico] 2023-24.
        \item[Grado] Doble Grado en Ingeniería Informática y Matemáticas.
        \item[Grupo] Único.
        \item[Profesor] Rafael Ortega Ríos.
        \item[Descripción] Primer parcial.
        \item[Fecha] 31 de octubre de 2023.
        %\item[Duración] 60 minutos.
    
    \end{description}
    \newpage
    

    \begin{ejercicio}
        Pruebe que la siguiente ecuación define una única función implícita $x: \bb{R} \to \bb{R}$, $t \mapsto x(t)$:
        \begin{equation*}
            e^x + x^3 + t = 0
        \end{equation*}
        Pruebe además que la función $x(t)$ es decreciente.\\

        Para que la ecuación anterior defina una única función implícita, hemos de ver que, para cada $t\in\bb{R}$, la ecuación $e^x + x^3 + t = 0$ tiene una única solución.
        Demostremos por tanto la existencia y unicidad de la solución de la ecuación anterior.
        \begin{description}
            \item[Existencia] 
            Definimos:
            \Func{f_t}{\bb{R}}{\bb{R}}{x}{e^x + x^3 + t}

            Tenemos que:
            \begin{equation*}
                \lim_{x\to-\infty} f_t(x) = -\infty, \quad \lim_{x\to+\infty} f_t(x) = +\infty
            \end{equation*}

            Como $f_t$ es continua, por el Teorema de Bolzano tenemos que existe $x_t\in\bb{R}$ tal que $f_t(x_t) = 0$. Por tanto, la ecuación $e^x + x^3 + t = 0$ tiene solución para cada $t\in\bb{R}$.

            \item[Unicidad]
            
            Como $f_t$ es derivable, tenemos que:
            \begin{equation*}
                f_t'(x) = e^x + 3x^2 > 0, \quad \forall x\in\bb{R}
            \end{equation*}

            Por tanto, $f_t$ es estrictamente creciente, por lo que es inyectiva. Por tanto, la ecuación $e^x + x^3 + t = 0$ tiene una única solución $x_t$ para cada $t\in\bb{R}$.
        \end{description}
            
        Sea entonces la función implícita la siguiente:
        \Func{x}{\bb{R}}{\bb{R}}{t}{x_t}

        Para probar que $x(t)$ es decreciente, podemos hacerlo de dos formas:
        \begin{description}
            \item[Opción 1] Sea $t_1, t_2\in\bb{R}$ tales que $t_1 < t_2$. Hemos de probar que $x(t_1) > x(t_2)$.
            
            Como $x(t_1)$ y $x(t_2)$ son las soluciones de las ecuaciones $e^x + x^3 + t_1 = 0$ y $e^x + x^3 + t_2 = 0$, respectivamente, tenemos que:
            \begin{equation*}
                e^{x(t_1)} + (x(t_1))^3 + t_1 = 0, \quad e^{x(t_2)} + (x(t_2))^3 + t_2 = 0
            \end{equation*}

            Restando ambas ecuaciones, obtenemos:
            \begin{equation*}
                e^{x(t_1)} + (x(t_1))^3 - e^{x(t_2)} - (x(t_2))^3 + t_1 - t_2 = 0
            \end{equation*}

            Como $t_1-t_2 < 0$, tenemos que:
            \begin{equation*}
                e^{x(t_1)} + (x(t_1))^3 - e^{x(t_2)} - (x(t_2))^3 > 0 \Longrightarrow
                e^{x(t_1)} + (x(t_1))^3 > e^{x(t_2)} + (x(t_2))^3
            \end{equation*}

            Definimos ahora la siguiente función:
            \Func{g}{\bb{R}}{\bb{R}}{x}{e^x + x^3}

            Como $g$ es suma de una función estrictamente creciente y otra función creciente, tenemos que $g$ es estrictamente creciente. Por tanto, usando $g$, tenemos que:
            \begin{equation*}
                g(x(t_1)) > g(x(t_2))
            \end{equation*}

            Como $g$ es estrictamente creciente, esto solo será posible si $x(t_1) > x(t_2)$, por lo que hemos probado que $x(t)$ es decreciente.

            \item[Opción 2] Aplicar el Teorema de la Función Implícita. Para ello, definimos:
            \Func{F}{\bb{R}^2}{\bb{R}}{(t,x)}{e^x + x^3 + t}

            Calculamos las derivadas parciales:
            \begin{equation*}
                \dfrac{\partial F}{\partial x}(t,x) = e^x + 3x^2, \quad \dfrac{\partial F}{\partial t}(t,x) = 1, \quad \forall (t,x)\in\bb{R}^2
            \end{equation*}

            Por tanto, $F\in C^1(\bb{R}^2)$. Además, tenemos $\dfrac{\partial F}{\partial x} (t, x) > 0$ para todo $(t, x)\in\bb{R}^2$ (en particular, no se anula). Por tanto, podemos aplicar el Teorema de la Función Implícita en cualquier punto $(t_0, x_0)$ tal que $F(t_0, x_0) = 0$,
            obteniendo así que la función implícita $x(t)$ es derivable. Aplicaremos este teorema a los puntos de la forma $(t,x_t)=(t,x(t))$. Calculamos su derivada mediante derivación implícita:
            \begin{equation*}
                \dfrac{\partial F}{\partial t}(t, x(t)) + \dfrac{\partial F}{\partial x}(t, x(t))\cdot x'(t) = 0
            \end{equation*}

            Despejando, tenemos que:
            \begin{equation*}
                x'(t) = -\dfrac{\dfrac{\partial F}{\partial t}(t, x(t))}{\dfrac{\partial F}{\partial x}(t, x(t))} = -\dfrac{1}{e^{x(t)} + 3(x(t))^2} < 0 \quad \forall t\in\bb{R}
            \end{equation*}

            Por tanto, hemos probado que $x(t)$ es decreciente.
        \end{description}
    \end{ejercicio}


\begin{ejercicio}
    Se considera la siguiente función:
    \Func{F}{]0, +\infty[}{\bb{R}}{t}{\displaystyle \int_0^{\sqrt{t}} e^{s^2} ~ds}

    ¿Es $F$ de clase $C^1$? En caso afirmativo, calcula la derivada.\\

    Definimos las siguientes funciones:
    \Func{\varphi}{\bb{R}^+}{\bb{R}^+}{t}{\sqrt{t}}
    \Func{\psi}{\bb{R}^+}{\bb{R}}{x}{\displaystyle \int_0^xe^{s^2} ~ds}

    Tenemos que $\varphi\in C^1(\bb{R}^+)$ de forma directa, y $\psi\in C^1(\bb{R}^+)$ por el Teoremema Fundamental del Cálculo. Sus respectivas derivadas son:
    \begin{align*}
        \varphi'(t) &= \dfrac{1}{2\sqrt{t}}\\
        \psi'(x) &= e^{x^2}
    \end{align*}
    
    
    Por tanto, como $F = \psi\circ\varphi$, tenemos que $F\in C^1(\bb{R}^+)$. Para calcular su derivada, aplicamos la regla de la cadena:
    \begin{equation*}
        F'(t) = \psi'(\varphi(t))\cdot\varphi'(t) = e^{(\sqrt{t})^2}\cdot\dfrac{1}{2\sqrt{t}} = \dfrac{e^t}{2\sqrt{t}} \quad \forall t\in\bb{R}^+
    \end{equation*}
\end{ejercicio}


\begin{ejercicio}
    Encuentra la solución del problema de valores iniciales siguiente:
    \begin{equation*}
        \dot{x} = \left(\dfrac{x}{t}\right)^3 + \dfrac{x}{t} - 1, \quad x(1) = 1
    \end{equation*}
    ¿En qué intervalo está definida?\\

    Tenemos que se trata de una ecuación homogénea definida en $D=\bb{R}^+\times \bb{R}$. Definimos el siguiente cambio de variable:
    \Func{\varphi=(\varphi_1, \varphi_2)}{\bb{R}^+\times \bb{R}}{D_1}{(t, x)}{(t,y)=\left(t, \dfrac{x}{t}\right)}

    Veamos el codomidio de $\varphi$:
    \begin{align*}
        D_1=\varphi(D) &= \left\{(t, y)\in\bb{R}^2 \mid (t,yt)\in D\right\}=D
    \end{align*}    

    Además, como podemos despejar cada componente de forma única, tenemos que $\varphi^{-1}$ es:
    \Func{\varphi^{-1}}{D_1}{D}{(t, y)}{(t,x)=\left(t, yt\right)}

    En primer lugar, tenemos que $\varphi,\varphi^{-1}$ son biyectivas. Además, ambas son de clase $C^1$, por lo que $\varphi$ es un difeomorfismo.
    Además, es admisible por no cambiar la primera variable. La ecuación transformada es:
    \begin{equation*}
        \dfrac{dy}{dt} = \dfrac{dy}{dx}=-\dfrac{x}{t^2} +\dfrac{x'}{t} = -\frac{y}{t} + \dfrac{y^3+y-1}{t} = \dfrac{y^3-1}{t} \text{ con dominio } D_1
    \end{equation*}

    Tenemos que se trata de una ecuación de variables separadas. Veamos en qué puntos se anula la función dependiente de $y$:
    \begin{equation*}
        y^3-1=0 \Longrightarrow y=1
    \end{equation*}

    Por tanto, tenemos que una solución de la ecuación diferencial dada es:
    \begin{equation*}
        y(t)=1\qquad \forall t\in \bb{R}^+
    \end{equation*}

    Deshaciedo el cambio de variable, una solución de la ecuación original es:
    \begin{equation*}
        x(t)=y(t)t = t \qquad \forall t\in \bb{R}^+
    \end{equation*}

    Tenemos que dicha solución satisfece $x(1)=1$, luego es la solución buscada. Está definida en $\bb{R}^+$, ya que así se tiene que $x\in C^1(\bb{R}^+)$ y $(t,x(t))\in D$ para todo $t\in \bb{R}^+$.
\end{ejercicio}


\begin{ejercicio}
    Demuestra que las fórmulas
    \begin{equation*}
        s = -e^t, \quad y = (t^2 + 1)x
    \end{equation*}
    definen un difeomorfismo que va de $D = \bb{R}^2$ a un dominio $\hat{D}$ que se especificará.
    Prueba que se trata de un cambio admisible para la ecuación $x' = x + t$ y encuentra la ecuación transformada.\\

    Definimos el siguiente cambio de variable:
    \Func{\varphi=(\varphi_1, \varphi_2)}{\bb{R}^2}{\wh{D}}{(t, x)}{(s, y)=\left(-e^t, (t^2+1)x\right)}

    Además, como podemos despejar cada componente de forma única, tenemos que $\varphi^{-1}$ es:
    \Func{\varphi^{-1}}{\wh{D}}{D}{(s, y)}{(t, x)=\left(\ln(-s), \dfrac{y}{t^2+1}\right)=\left(\ln(-s), \dfrac{y}{\ln^2(-s)+1}\right)}

    Veamos el codominio de $\varphi$:
    \begin{align*}
        \wh{D}=\varphi(D) &= \left\{(s, y)\in\bb{R}^2 \mid \left(\ln(-s), \dfrac{y}{\ln^2(-s)+1}\right)\in D\right\}
    \end{align*}

    En primer lugar, necesitamos que la inversa esté bien definida, luego $s<0$. Además, el denominador de la segunda componente de esta no se anula nunca. Como $D=\bb{R}^2$, tenemos que $\wh{D}=\bb{R}^- \times \bb{R}$.\\

    Una vez visto que $\varphi,\varphi^{-1}$ son biyectivas, falta por ver que son de clase $C^1$. La primera componente de ambas es una función elemental de clase $C^1$, y la segunda componente es un producto o cociente de funciones de clase $C^1$, que en definitiva es de clase $C^1$.
    Por tanto, como ambas componentes de $\varphi, \varphi^{-1}$ son de clase $C^1$, tenemos que $\varphi, \varphi^{-1}$ son de clase $C^1$, y por tanto $\varphi$ es un difeomorfismo.\\

    Veamos ahora que se trata de un cambio admisible para la ecuación $x' = x + t$. Tenemos que:
    \begin{equation*}
        \dfrac{\partial \varphi_1}{\partial t} + \dfrac{\partial \varphi_1}{\partial x} x'=-e^{-t}+0\cdot x'=-e^{-t}\neq 0 \qquad \forall (t,x)\in D
    \end{equation*}

    Por tanto, se trata de un cambio admisible. La ecuación transformada es:
    \begin{equation*}
        \dfrac{dy}{ds} = \dfrac{dy}{dt}\cdot \dfrac{dt}{ds} = \dfrac{2tx+(t^2+1)x'}{-e^t} = \dfrac{2tx+(t^2+1)(x+t)}{-e^t}
    \end{equation*}

    Aplicamos ahora $\varphi,\varphi^{-1}$ para que la ecuación esté en función de $y$ y $s$:
    \begin{equation*}
        y'= \dfrac{2tx+y+(t^2+1)t}{s} = \dfrac{2\ln(-s)\cdot \dfrac{y}{\ln^2(-s)+1}+y+(\ln^2(-s)+1)\ln(-s)}{s}
    \end{equation*}

    Por tanto, la ecuación transformada es:
    \begin{equation*}
        y' = \dfrac{2\ln(-s)\cdot \dfrac{y}{\ln^2(-s)+1}+y+(\ln^2(-s)+1)\ln(-s)}{s}\qquad \text{con dominio }\wh{D}
    \end{equation*}
\end{ejercicio}

\begin{ejercicio}
    Se considera la transformación en el plano
    \begin{equation*}
        \psi(\theta, r) = (t, x), \quad t = r\cos \theta, \quad x = r\sen \theta, \quad (\theta, r) \in \wh{\Omega} = \left]\nicefrac{-\pi}{2}, \nicefrac{\pi}{2}\right[ \times \left]0, +\infty\right[
    \end{equation*}

    Determina $\Omega = \psi(\wh{\Omega})$ y prueba que $\psi$ es un difeomorfismo de $\wh{\Omega}$ a $\Omega$. Dada una ecuación $\dfrac{dx}{dt} = f(t, x)$ con $f : \Omega \to \bb{R}$, ¿bajo qué condiciones se puede asegurar que el difeomorfismo $\varphi = \psi^{-1}$ es admisible?\\

    Vemos directamente que $\psi$ es el cambio de coordenadas de polares a cartesianas en el semiplano $x>0$, es decir, primer y cuarto cuadrante.
    Busquemos en primer lugar la inversa de $\psi$. Buscamos despejar $\theta, r$ en función de $t, x$.
    \begin{itemize}
        \item Dividimos $x$ entre $t$, ya que el coseno no se anula en el dominio de defición dado. Tenemos que:
        \begin{equation*}
            \dfrac{x}{t} = \dfrac{r\sen \theta}{r\cos \theta} = \tan \theta
            \Longrightarrow \theta = \arctan \left(\dfrac{x}{t}\right)
        \end{equation*}
        donde en la última implicación hemos usado que $\theta\in\left]\nicefrac{-\pi}{2}, \nicefrac{\pi}{2}\right[$.
        Concluimos además que $t\neq 0$, y como $t=r\cos \theta$, tenemos que $t>0$.

        \item Despejamos ahora $r$. Sumando $x^2$ + $t^2$, tenemos:
        \begin{equation*}
            x^2+t^2 = r^2(\sen^2 \theta + \cos^2 \theta) = r^2 \Longrightarrow r = \sqrt{x^2+t^2}
        \end{equation*}
        donde hemos usado que $r>0$.
    \end{itemize}

    Por tanto, como hemos podido despejar cada componente de forma única, tenemos que $\psi^{-1}=\varphi$ es:
    \Func{\psi^{-1}=\varphi=(\varphi_1, \varphi_2)}{\Omega}{\wh{\Omega}}{(t, x)}{(\theta, r)=\left(\arctan \left(\dfrac{x}{t}\right), \sqrt{x^2+t^2}\right)}

    Estudiemos ahora el conjunto $\wh{\Omega}$. Tenemos que:
    \begin{align*}
        \Omega = \psi(\wh{\Omega}) &= \left\{(t, x)\in\bb{R}+\times \bb{R} \mid \left(\arctan \left(\dfrac{x}{t}\right), \sqrt{x^2+t^2}\right)\in \wh{\Omega}\right\}
        =\\&= \left\{(t, x)\in\bb{R}^+\times \bb{R} \mid \arctan \left(\dfrac{x}{t}\right)\in\left]\nicefrac{-\pi}{2}, \nicefrac{\pi}{2}\right[, \sqrt{x^2+t^2}>0\right\}
        =\\&= \bb{R}^+\times \bb{R}
    \end{align*}

    Ya sabemos que $\psi, \psi^{-1}$ son biyectivas. Además, $\psi,\psi^{-1} \in C^1$ de forma directa, puesto que son composición de funciones de clase $C^1$. Por tanto, $\psi,\psi^{-1}$ son difeomorfismos.\\

    Para que $\varphi$ sea admisible para dicha ecuación, necesitamos que se cumpla:
    \begin{align*}
        0&\neq \dfrac{\partial \varphi_1}{\partial t} + \dfrac{\partial \varphi_1}{\partial x} x'
        = \dfrac{1}{1+\left(\nicefrac{x}{t}\right)^2}\cdot \left(-\dfrac{x}{t^2}\right) + \dfrac{1}{1+\left(\nicefrac{x}{t}\right)^2}\cdot \dfrac{1}{t}\cdot f(t, x)
        =\\&= \dfrac{1}{1+\left(\nicefrac{x}{t}\right)^2}\left(-\dfrac{x}{t^2} + \dfrac{f(t, x)}{t}\right)
        \Longleftrightarrow
        \dfrac{-x}{t}+f(t, x)\neq 0
        \Longleftrightarrow
        f(t, x)\neq \dfrac{x}{t} \qquad \forall (t, x)\in \Omega
    \end{align*}

    Por tanto, la condición de admisibilidad equivale a:
    \begin{equation*}
        f(t, x)\neq \dfrac{x}{t} \qquad \forall (t, x)\in \bb{R}^+\times \bb{R}
    \end{equation*}
\end{ejercicio}

    
\end{document}